% GENERATED FILE. Edit canonical lessons, then run npm run generate:books.
% canonical-source-hash: fnv1a64:8c5ffb0b0752113c
% canonical-lessons: ES-R01-no-entiendo, ES-R02-mas-despacio, ES-R03-repaso-reparar

\chapter{When You Don't Understand}
\label{ch:repair}
\hlchaptermodality{14}

\begin{chapteropening}
\textbf{By the end of this chapter:} \emph{I can say that I did not understand, and ask for it more slowly.}

The two sentences that rescue a conversation, and the thanks that closes it.
\end{chapteropening}

\section[no entiendo]{no entiendo — the most useful sentence a beginner owns}

\label{lesson:ES-R01-no-entiendo}



\begin{quote}
Everything so far assumed you understood what was said to you. This chapter is for when you did not.
\end{quote}

\subsection*{What to know first}
\textbf{No entiendo.} — \textquotedblleft{}I don't understand.\textquotedblright{}

Say it whole. You already have \textbf{no}; the rest comes with it for now, and the piece you are not yet meant to take apart will come apart later.

This is not an admission of failure. It is the sentence that keeps a conversation alive — the alternative is nodding at something you did not catch.

\begin{cousinweb}
The word hiding inside \emph{entiendo} comes from Latin \textbf{intendere} — literally \emph{to stretch toward}. To understand something was to lean into it.

English took the same word and kept it in several places: \textbf{intend} (to stretch your mind toward a thing), \textbf{intense}, and — through the same family — \textbf{attention}, which is what you are stretching when you listen hard.

So \emph{no entiendo} is, underneath, \textquotedblleft{}I am not managing to stretch that far yet.\textquotedblright{} Which is exactly true, and much kinder than \textquotedblleft{}I failed.\textquotedblright{}
\end{cousinweb}

\subsection*{Your turn}
Say these aloud:

\begin{itemize}
\raggedright
  \item \textquotedblleft{}No entiendo.\textquotedblright{} — calmly, not apologetically
  \item it again, to someone speaking too fast
\end{itemize}

\begin{tcolorbox}[breakable,colback=teal!4,colframe=teal!35!black,title={Before you move on}]
Say it once from memory. Which English word shares its root?
\end{tcolorbox}
\section[más despacio]{más despacio — the request that makes the last sentence solvable}

\label{lesson:ES-R02-mas-despacio}



\begin{quote}
Everything so far assumed you understood what was said to you. This chapter is for when you did not.
\end{quote}

\subsection*{What to know first}
\textbf{Más despacio, por favor.} — \textquotedblleft{}More slowly, please.\textquotedblright{}

\emph{No entiendo} stops the conversation. This restarts it. Together they are the whole repair kit: say what went wrong, then say what would fix it.

\begin{cousinweb}
\emph{Despacio} comes from \textbf{de spatio} — \textquotedblleft{}of space.\textquotedblright{} Latin \textbf{spatium} is English \textbf{space}, and also \textbf{spatial}.

Doing something \emph{de spatio} meant doing it with room around it, unhurried. The sense drifted from space into time, which is a journey English made too: we still say we need \emph{space} when we mean we need \emph{time}.
\end{cousinweb}

\subsection*{Your turn}
Say these aloud:

\begin{itemize}
\raggedright
  \item \textquotedblleft{}Más despacio.\textquotedblright{} — calmly, not apologetically
  \item it again, to someone speaking too fast
\end{itemize}

\begin{tcolorbox}[breakable,colback=teal!4,colframe=teal!35!black,title={Before you move on}]
Say it once from memory. Which English word shares its root?
\end{tcolorbox}
\section[Practice]{When you don't understand — putting it together}

\label{lesson:ES-R03-repaso-reparar}



\begin{quote}
Two sentences. Between them they rescue almost any conversation.
\end{quote}

\subsection*{The exchange}
\begin{quote}
— \emph{(something too fast to follow)} — \textbf{No entiendo. Más despacio, por favor.} — \emph{(again, slower)} — \textbf{Gracias.}
\end{quote}

\begin{grammarlens}[title={What you've built}]
Until now, a sentence you missed was the end of the conversation. It is not any more. You can say what went wrong and what would fix it, and then thank someone for doing it — which is the difference between speaking a language badly and not speaking it at all.
\end{grammarlens}

\subsection*{Your turn}
Say these aloud:

\begin{itemize}
\raggedright
  \item both sentences, one after the other, without a pause between them
  \item them again, then \textquotedblleft{}Gracias.\textquotedblright{} — the whole repair, start to finish
\end{itemize}

\emph{Twice through:}

\begin{tcolorbox}[breakable,colback=teal!4,colframe=teal!35!black,title={Before you move on}]
What do you say when you did not follow? And what do you ask for next? Which English words hide inside \emph{entiendo} and \emph{despacio}? (\emph{Intend, intense, attention} — and \emph{space}.)
\end{tcolorbox}
